% Response_to_Reviewers.tex
\documentclass[11pt,a4paper]{article}

\usepackage[utf8]{inputenc}
\usepackage[T1]{fontenc}
\usepackage{lmodern}
\usepackage{microtype}

\usepackage{amsmath,amssymb}
\usepackage{geometry}
\usepackage{enumitem}
\usepackage{needspace}
\usepackage{parskip}
\usepackage{hyperref}
\usepackage{xurl}
\usepackage{xcolor}

\geometry{
  top=2.5cm,
  bottom=2.5cm,
  left=2.7cm,
  right=2.7cm
}

\hypersetup{
  colorlinks=true,
  linkcolor=black,
  urlcolor=blue,
  citecolor=black
}

\setlength{\parindent}{0pt}
\setlength{\parskip}{0.7em}
\clubpenalty=10000
\widowpenalty=10000
\displaywidowpenalty=10000
\let\originalsubsection\subsection
\renewcommand{\subsection}{\Needspace{9\baselineskip}\originalsubsection}
\let\originalsection\section
\renewcommand{\section}{\Needspace{12\baselineskip}\originalsection}

% Reviewer-comment formatting
\newenvironment{reviewcomment}
  {\begin{quote}\itshape}
  {\end{quote}}

\newcommand{\response}{\noindent\textbf{Response:} }

\title{\textbf{Response to the Editor and Reviewers}}
\author{}
\date{}

\begin{document}

\maketitle
Submission ID: \texttt{1e58c388-d7ec-4084-b792-df20a48e3e75}.
Locations refer to section titles and figure/table labels.

We thank the editor and reviewers for their careful reading. We have
clarified the horizon definitions, explained their practical use, and
revised the presentation of the four-domain results. We also checked the
stored evidence behind the numerical summaries and the alternative-baseline
analysis. The corrections described below distinguish changes to the
reporting from changes to experiments.

\textbf{Additional experiments carried out in this revision.} No new model
class, hyperparameter, split or evaluation protocol was introduced. We carried out the following
analyses, all on the same data and protocol as the paper:
\begin{enumerate}[leftmargin=*,itemsep=2pt]
\item \emph{Wind ARIMA check with a true $h$-step forecast.} The earlier check requested a one-step forecast at
every horizon. We reran it with the $h$-step forecast; the wind descriptors are unchanged (Section~Wind Domain Results).
\item \emph{ARIMA order diagnostic.} An AIC search ($p,q\le5$, $d=0$) on the training windows of 20 evenly spaced
evaluated origins per domain (Appendix~A.2). The reported ARIMA check was not refitted with these orders.
\item \emph{RMSE recomputation.} The four LightGBM experiments were re-executed unchanged while also recording RMSE.
Their MAE reproduces the reported MAE exactly, and Appendix~A.1 now uses these per-origin errors.
\item \emph{Seasonal-persistence sensitivity.} Re-executed on the current inputs; all eight reported
descriptor tuples are reproduced exactly (Table~4).
\end{enumerate}


% ============================================================
\section*{Editor}
% ============================================================

\begin{reviewcomment}
Please report results precisely, remove overstated conclusions, and explain
limitations completely.
\end{reviewcomment}

\response
We corrected a mismatch between the PM$_{2.5}$ text and its plotted
LightGBM curve. That curve gives $H^*(\mathrm{strict})=22$ and $[27,48]$.
The former value of 13 and $[36,48]$ came from a separate moving-average
result, not the reported LightGBM curve. We corrected the text, caption,
and Table~3 without changing the forecasts or the figure data.

In ``Models'', the daily-load description now follows the producer:
expanding training, evaluation beginning after 365 observations, and a
weekly index-phase feature rather than a timestamp-derived day-of-week
label. We removed claims that the observed profiles establish intrinsic
domain limits or robustness across forecasting models. ``Limitations''
now separates positive average skill from evidence of operational benefit
and makes the baseline, metric, support, model, and horizon caps explicit.
The discussion of dynamic horizon selection remains a potential use,
not a validated controller.

The archived four-domain seasonal-baseline analysis is reported in
``Alternative-baseline sensitivity'' and Table~4. Its stored comparisons
were checked on exact common support, and the analysis was re-executed on the
current inputs, reproducing all reported values exactly.
The metric-sensitivity appendix (Appendix~A.1) was recomputed from per-origin
squared errors of the same LightGBM forecasts. The rerun reproduces the reported
MAE exactly in all four domains, so the MAE and RMSE descriptors share forecasts
and support. This corrected the PM$_{2.5}$ RMSE row (48/23, $[26,48]$) that had
been inherited from a different, moving-average experiment; the other three rows
are unchanged (wind 48/48, traffic 72/63, load 3/3).


% ============================================================
\section*{Reviewer 1}
% ============================================================

\subsection*{Comment 1}

\begin{reviewcomment}
In the operational definitions, use a colon before the explanations of the
relaxed and strict horizons.
\end{reviewcomment}

\response
We replaced the full stop with a colon after both labels in
``Operational Predictability Horizon'' (Section~3.1; the reviewed PDF numbers the "Operational definitions" paragraph as Section~3.1.1).


\subsection*{Comments 2--3}

\begin{reviewcomment}
Equation (5) and the surrounding text should explicitly state that the relaxed
horizon permits non-positive skill at intermediate steps.
\end{reviewcomment}

\response
We now define
\[
\mathcal{H}^{+}
=
\{h\in\mathcal{H}:\mathrm{Skill}(h)>0\}
\]
and
\[
H^*(\mathrm{relax})
=
\max\left(\mathcal{H}^{+}\cup\{0\}\right)
\]
in Section~3.1. Equation~(5) also defines the relaxed span $\mathcal{R}$
and the gap set $\mathcal{G}=\mathcal{R}\setminus\mathcal{H}^{+}$,
whose elements have non-positive skill. This set may be non-empty.

The text explicitly notes that intermediate horizons with non-positive skill
may fall within the relaxed reach, but are not themselves interpreted as
skilful. We verified this behaviour using executable unit tests covering
intermediate skill deficits and delayed onset.


\subsection*{Comment 4}

\begin{reviewcomment}
Maintain units throughout Table 2.
\end{reviewcomment}

\response
Table~2 now specifies PM$_{2.5}$ in $\mu\mathrm{g\,m^{-3}}$, wind speed in
$\mathrm{m\,s^{-1}}$, and traffic speed in mph. For electric load, we clarify that the
stored experimental target is the sum of 15-minute kW readings and that
division by four expresses the aggregate daily demand in kWh, following the source
metadata and the aggregation script.


\subsection*{Comment 5}

\begin{reviewcomment}
Figures 1--3 are not correctly referenced in the text.
\end{reviewcomment}

\response
We have added explicit textual references to the conceptual $H^*$ schematic
(Fig.~1), the rolling-origin validation protocol (Fig.~2), and the
PM$_{2.5}$ skill curve (Fig.~3).


\subsection*{Comment 6}

\begin{reviewcomment}
Evaluate baseline choices beyond persistence.
\end{reviewcomment}

\response
We include the archived seasonal-persistence analysis, using a fixed
24-hour period for the hourly domains and a 7-day period for daily load.
For each domain, we verified equality of forecast origin, target timestamp,
horizon, and $y_{\mathrm{true}}$ between the stored LightGBM, persistence,
and seasonal-persistence records. Pairwise alignment loses no records.
The analysis was re-executed on the current inputs and reproduces all eight
reported descriptor tuples exactly.

For PM$_{2.5}$, the persistence comparison on this exact common support yields
$H^*(\mathrm{strict})=22$ with longest contiguous positive interval $[27,48]$,
whereas seasonal persistence yields $H^*(\mathrm{strict})=24$ with $[25,48]$.

The primary LightGBM curve also gives $H^*(\mathrm{strict})=22$ and
$[27,48]$, after the reporting correction explained above. The previous
13-hour value is not explained as a support effect. Equal descriptors
do not imply equal samples or MAEs: the seasonal-baseline comparison
uses only its matched-support rows.

Under the same matched-support analysis, load remains at
$H^*(\mathrm{strict})=1$, wind changes from 48 to 30, and traffic changes from
7 to 14. $H^*(\mathrm{relax})$ remains unchanged across all four domains. We
note that the hourly endpoints are exact multiples of the seasonal period,
where seasonal and standard persistence coincide; the unchanged relaxed
endpoint is therefore partly structural and is not interpreted as evidence of
baseline invariance.


\subsection*{Comment 7}

\begin{reviewcomment}
Explain why ARIMA(2,0,0) was selected and whether ACF was examined.
\end{reviewcomment}

\response
We agree that the original manuscript did not document this choice. The order
$(2,0,0)$ was fixed a priori as a low-order autoregression; an audit of our
repository found no record that it was selected via ACF/PACF, information
criteria, grid search, auto-ARIMA or validation splits, and we do not claim
otherwise. The ARIMA check is now described in Section~Models as a
model-substitution illustration, not a tuned competitor.

To answer the question directly, we examined the data. On the training windows
of 20 evenly spaced evaluated origins per domain, an AIC search
($p,q\le5$, $d=0$, with constant) selects different orders at different origins
(10 distinct best orders in wind, 8 in traffic), and $(2,0,0)$ lies within
$\Delta$AIC$\,\le 2$ of the optimum in 10 of 20 wind windows but only 4 of 20
traffic windows (Appendix~A.2). We report this as a limitation of the fixed
order, not as a justification. The ACF and PACF of the last training window of each domain
are shown in a new figure in Appendix~A.2; they were inspected after the order was fixed. The PACF is dominated by lag~1,
and the traffic ACF has a daily peak at lag~24 that a non-seasonal AR(2) cannot represent.

The audit also found that the historical wind check requested a one-step
forecast at every horizon instead of the stated $h$-step forecast. We reran it
with a true $h$-step forecast, otherwise identical: the wind descriptors are
unchanged ($H^*(\mathrm{relax})=H^*(\mathrm{strict})=48$, $[1,48]$, no sign
changes) while skill amplitude differs. Traffic used an $h$-step forecast and its
numbers are unchanged ($H^*(\mathrm{strict})=6$, $[38,43]$, 13 sign changes).

The separate seasonal-baseline analysis tests sensitivity to the reference
baseline, not robustness across forecasting models, and is not used to justify
the ARIMA order.


\subsection*{Comments 8--9}

\begin{reviewcomment}
Explain the practical utility of relaxed and strict $H^*$ and what they add
beyond fixed-horizon summaries.
\end{reviewcomment}

\response
A new ``Operational interpretation and actionability'' subsection
(Section~3.2) defines $H^*(\mathrm{relax})$ as the outer horizon at
which positive baseline-relative skill is observed, and $H^*(\mathrm{strict})$,
with its interval location, as the longest contiguous positive interval.

We explain that a few fixed-horizon errors do not locate positive-skill
intervals or show the gaps between them. The PM$_{2.5}$ and traffic
examples illustrate why interval locations are needed in addition to
lengths. We also state that positive average skill alone does not establish
deployment value: its magnitude, uncertainty, decision loss, and cost matter.


\subsection*{Comment 10}

\begin{reviewcomment}
Could $H^*$ dynamically decide which horizons to forecast, and could this
create feedback or overfitting?
\end{reviewcomment}

\response
We discuss this possibility in ``Discussion''. In an
operational setting, $H^*$ could inform a decision policy estimated on
historical validation data, frozen for prospective deployment, and subsequently
updated using only realised historical errors.

We explicitly caution that estimating $H^*$, selecting horizons, and
evaluating performance on the same errors introduces adaptive selection bias.
A dynamic controller would require prospective, prequential, or nested
validation, and is not claimed to be validated here.


% ============================================================
\section*{Reviewer 2}
% ============================================================

\subsection*{Comment 1}

\begin{reviewcomment}
Explain the actionability of the findings more explicitly.
\end{reviewcomment}

\response
We added the dedicated operational interpretation subsection noted above
(Section~3.2), distinguishing the outer positive reach from the longest
contiguous positive interval and restricting dynamic deployment claims to
future, separately validated policies.


\subsection*{Comment 2}

\begin{reviewcomment}
Figures 1 and 2 are not referenced correctly.
\end{reviewcomment}

\response
Both figures are now explicitly introduced and discussed in the main text
(Figs.~1--2).


\subsection*{Comment 3}

\begin{reviewcomment}
In Figure 1, the curve overlaps text or annotations.
\end{reviewcomment}

\response
The revised Figure~1 uses dedicated callout boxes and separated label
placements. This is a layout correction, not a change to the horizon
definitions or empirical results.


\subsection*{Comment 4}

\begin{reviewcomment}
In Figure 2, text does not fit inside the boxes.
\end{reviewcomment}

\response
The revised Figure~2 has larger boxes and wrapped text.


\subsection*{Comment 5}

\begin{reviewcomment}
If possible, evaluate alternative baselines beyond persistence.
\end{reviewcomment}

\response
We report the archived seasonal-persistence analysis
detailed in our response to Reviewer~1 (Comment~6), enforcing exact common
support and reporting all outcomes regardless of direction. Table~4 reports
all four domains across both baselines, documenting common-origin sample
counts, $H^*(\mathrm{relax})$, $H^*(\mathrm{strict})$, and interval locations.


\subsection*{Comment 6}

\begin{reviewcomment}
Please cache all the data in a site you control and make the central
data location clear in the paper.
\end{reviewcomment}

\response
We deposited the exact input files of all four experiments in the project
repository, with SHA-256 checksums and provenance notes, under the fixed tag
\texttt{data-cache-v1}:
\url{https://github.com/EduIntellect/paper2H/tree/data-cache-v1/release/data_cache_four_domain}.
The ``Data Availability'' section now points to this location. PM$_{2.5}$ and electric
load are redistributed under CC BY 4.0 and wind under NREL's open-data terms. The traffic file is a
single-sensor hourly aggregate of METR-LA; its underlying measurements come from Caltrans PeMS, and
we state in the repository that we did not obtain separate confirmation of redistribution terms and will
remove the file on request.


\subsection*{Comment 7}

\begin{reviewcomment}
Explain the reason for ARIMA(2,0,0).
\end{reviewcomment}

\response
As noted in our response to Reviewer~1 (Comment~7), the order $(2,0,0)$ was fixed
a priori as a low-order autoregression and was not selected with ACF/PACF or
information criteria. We now say so in Section~Models, report an a posteriori
AIC-based diagnostic in Appendix~A.2 (the best order varies across origins, and
$(2,0,0)$ is further from the optimum in traffic than in wind), and present the
check as an illustration of descriptor sensitivity, not a tuned model. The
wind check was rerun with a true $h$-step forecast; its descriptors are
unchanged.


\end{document}
